% !TeX root = tesis.tex

\subsection{De protocolos interactivos a no interactivos}

Los protocolos de argumentos de conocimiento suelen ser interactivos entre el prover y el verifier. Sin ir más lejos, el argumento sucinto desarrollado en la sección anterior requiere de que el verifier envíe un valor aleatorio como challenge al prover.

Sin embargo, en entornos como blockchains, donde la verificación debe ocurrir en un único mensaje publicado onchain, la interactividad es un obstáculo. En 1986, Amos Fiat y Adi Shamir propusieron una transformación que convierte muchos de estos protocolos interactivos en protocolos no interactivos (\NIZK) sin pérdida significativa del valor de seguridad, cuando se modela la función de desafío como un oráculo aleatorio \cite{Fiat, Bellare1993}.

La idea central es la siguiente: en lugar de que el \textit{verifier} genere un desafío aleatorio y se lo envíe al \textit{prover}, el \textit{prover} calcula el desafío como el resultado de aplicar una función de hash criptográfica a la transcripción parcial del protocolo y utiliza ese resultado como si fuera el desafío.

Para que esta transformación sea segura, se necesita trabajar en el \ROM, de modo que la función de hash $\Hash$ se comporte como una función aleatoria. Este enfoque es central para la mayoría de construcciones de \zkSNARK modernas: aunque no existen funciones hash verdaderamente aleatorias en la práctica, asumiendo que una función criptográfica modelada como oráculo aleatorio permite demostrar seguridad de las construcciones de manera formal y luego instanciarlas con hashes como \SHATwoFiveSix, \BLAKETwo o \PoseidonHash bajo heurísticas de seguridad razonables. Para formalizar esta transformación, primero planteamos formalmente un protocolo interactivo público genérico de tres mensajes (lo cual puede trivialmente generalizarse a un número arbitrario de mensajes), conocido como \SigmaProtocol.

\begin{definition}[\SigmaProtocol]
	Decimos que $\Pi = (\Commit, \Open, \Verify)$ es un \SigmaProtocol para una relación $\Relation \subseteq \InstanceSpace \times \WitnessSpace$ con public inputs $x$ y witness $w$ tales que $(x,w) \in \Relation$, entre un prover $\Prover$ y un verifier $\Verifier$, si consiste de los siguientes tres mensajes:
	\begin{enumerate}[itemsep=0.1 cm]
		\item El prover le envía al verifier un compromiso $c = \Commit(x, w; r)$ ($r$ es la aleatoriedad interna del prover, fijada al inicio de la ejecución).
		\item El verifier responde con un desafío uniformemente aleatorio $z$.
		\item El prover le envía una respuesta $\pi = \Open(c, z, w, x; r)$.
	\end{enumerate}
	Una ejecución completa del $\Sigma$-protocol produce la tripla $(c, z, \pi)$. El verifier acepta la respuesta $\pi$ si $\Verify(x, c, z, \pi) = 1$.
\end{definition}

\begin{definition}[Transformación \FiatShamir]
	Sea $\Pi = (\Commit, \Open, \Verify)$ un \SigmaProtocol para una relación $\Relation \subseteq \InstanceSpace \times \WitnessSpace$. La \emph{transformación \FiatShamir} bajo una función de hash criptográfica $\Hash: \Strings{*} \to \Strings{n}$ es un protocolo no interactivo $\Pi_{FS} = (\Prove_{FS}, \Verify_{FS})$ definido como:
	\begin{itemize}[noitemsep]
		\item $\Prove_{FS}(x,w)$, ejecutado por el prover:
		\begin{itemize}[noitemsep]
			\item Se computa $c = \text{Commit}(x,w;r)$.
			\item Se computa un \textit{challenge} $\Challenge \in \Strings{n}$ como $\Challenge = \Hash(x \concat c)$, donde la concatenación $x \concat c$ representa una codificación no ambigua del \textit{transcript} parcial.
			\item Se computa la respuesta $\pi = \Open(c, z, w, x; r)$.
		\end{itemize}
		\item $\Verify_{FS}(x,c,\pi)$, ejecutado por el verifier:
		\begin{itemize}[noitemsep]
			\item Se recomputa el challenge $\Challenge = \Hash(x \concat c)$.
			\item Se acepta si y solo si $\Verify(x,c,z,\pi) = 1$.
		\end{itemize}
	\end{itemize}
\end{definition}

\subsection{Transformación \FiatShamir en \KZG}

Vamos a ver un ejemplo de la transformación de \FiatShamir. En particular, la vamos a aplicar al esquema de compromisos polinomiales \KZG. Esto nos permitirá ganar intuición sobre la transformación, lo que será necesario para la sección final del capítulo. Observemos que esto es para las etapas posteriores al trusted setup.

Sea una función hash criptográfica $\Hash: \Strings{*} \to \Fp$, modelada como oráculo aleatorio en el \ROM. La transformación consiste en reemplazar el desafío aleatorio $\Challenge \in \Fp$ por $\Challenge := \Hash(x \concat c)$, donde $c$ es el compromiso polinomial generado por el prover. En implementaciones reales, se computa $\Challenge :=  \Hash(x \concat \textnormal{encode}(c))$, para una codificación determinística y no ambigua.

El protocolo consiste de los algoritmos $\Prove_{FS}$ \ref{alg:kzg-fs-prove}, a ser ejecutado por el prover, y $\Verify_{FS}$ \ref{alg:kzg-fs-verify}, a ser ejecutado por el verifier.

\begin{algorithm}[ht]
	\caption{$\Prove_{FS}$ para el esquema \KZG}
	\label{alg:kzg-fs-prove}
	\begin{algorithmic}[1]
		\Require Polinomio $P \in \poly{\Fp}_{\leq d}$, \SRS $\{G_1, G_1\tau, \ldots, \tau^d G_1, G_2, G_2 \tau\}$.
		\Ensure $(c, y, \pi)$ es una prueba no interactiva que será aceptada por el algoritmo $\Verify_{FS}$.
		\State $c \gets P(\tau) G_1$ \Comment{cómputo del compromiso $c$}
		\State $\Challenge \gets \Hash(x \concat c)$ \Comment{cómputo del \textit{challenge} $\Challenge$}
		\State $y \gets P(\Challenge)$ \Comment{evaluación del polinomio $P$ en el \textit{challenge}}
		\State $Q(x) \gets \frac{P(x) - y}{x - \Challenge}$ \Comment{cómputo del polinomio cociente $Q$}
		\State $\pi \gets Q(\tau) G_1$ \Comment{cómputo del argumento de evaluación $\pi$}
		\State \Return $(c, y, \pi)$
	\end{algorithmic}
\end{algorithm}

\begin{algorithm}[ht]
	\caption{$\Verify_{FS}$ para el esquema \KZG}
	\label{alg:kzg-fs-verify}
	\begin{algorithmic}[1]
		\Require Compromiso polinomial $c \in \GOne$, evaluación $y \in \Fp$, argumento de evaluación $\pi \in \GOne$.
		\Ensure Si el algoritmo devuelve $1$, entonces el prover conoce un polinomio $P$ consistente con $c$ tal que $y = P(\Challenge)$, excepto por probabilidad negligible.
		\State $\Challenge \gets \Hash(x \concat c)$ \Comment{recómputo del \textit{challenge} $\Challenge$}
		\State \Return $\left\llbracket \pairing{c - yG_1}{G_2} = \pairing{\pi}{\tau - \Challenge) G_2} \right\rrbracket$
	\end{algorithmic}
\end{algorithm}

La transformación de \FiatShamir aplicada al protocolo de apertura de \KZG muestra un principio fundamental en el diseño de pruebas criptográficas modernas: la eliminación de la interacción mediante la derivación determinística de desafíos aleatorios a partir de información pública. Esta técnica permite eliminar la interactividad en protocolos sin sacrificar las demás propiedades deseables en un \SNARK.

En el caso del esquema \KZG, el desafío $\Challenge$ pasa de ser generado por el verificador a determinarse como el resultado de una función hash aplicada al compromiso. En contraste, como veremos en la siguiente sección, el esquema \Groth es inherentemente \NIZK: la aleatoriedad que en un \SigmaProtocol sería provista por el verificador se encuentra implícitamente embebida en el trusted setup. La transformación de \FiatShamir reaparecerá más adelante en este trabajo en otros contextos, donde protocolos interactivos deben ser compilados a versiones no interactivas.

\subsection{Limitaciones}

Aunque el esquema \FiatShamir es extremadamente útil y ampliamente utilizado para obtener pruebas no interactivas, no es una panacea sin riesgos. La seguridad de la transformación se demuestra únicamente en el \ROM, y no existe una garantía absoluta en el mundo real de que una instancia concreta del hash cumpla exactamente ese ideal. De hecho, múltiples trabajos han documentado varios problemas asociados a implementaciones mal hechas o a condiciones de seguridad insuficientes. La bibliografía es sumamente extensa al respecto; acá dejamos un pequeño resumen.

En particular, está demostrado que existen protocolos interactivos seguros para los cuales la transformación \FiatShamir falla para cualquier función hash real, lo que subraya que la seguridad depende críticamente de la idealización de la función de hash \cite{Goldwasser}.

En diversos trabajos recientes \cite{Dao2023, cryptoeprint:2024/1565, Khovratovich2025}, se analizan distintos tipos de errores en la aplicación de la transformación en sistemas reales, ilustrando ejemplos en los que la construcción incorrecta de los inputs al hash puede comprometer la seguridad de pruebas basadas en esta heurística. Los autores de estos trabajos demuestran que un uso débil o parcial de la transformación puede permitir ataques efectivos contra protocolos de \SNARK concretos, incluso pudiendo falsificar pruebas.

Finalmente, observamos que en el contexto de la computación cuántica, las transformaciones \FiatShamir con funciones de hash reales tienen limitaciones de seguridad bajo ataques que explotan superposiciones de consultas \cite{cryptoeprint:2024/590}.

Como lecciones principales para el desarrollo de nuestro trabajo, es fundamental especificar con precisión qué partes del transcript se incluyen en el hash, utilizar codificaciones no ambiguas y evitar reutilización de desafíos entre instancias distintas del protocolo.
